\begin{textsample}{NEG Dim 9 – global\_north\_2023\_12 – Score -2.00 – global\_north\_2023\_12/104573810.txt}  \label{ex:f9_neg_292}
\textbf{Newsletters} The New York Times is facing backlash after publishing an op-ed written by the Hamas-backed mayor of Gaza City describing the state of the Gaza Strip amid the war between Israel and the Hamas terrorist group that broke out in early October.
The news outlet published the op-ed on Sunday, written by Yahya Sarraj, mayor of Gaza, of which Hamas has been the de facto governing body for more than a decade.
The essay laments the actions of the Israeli military, particularly after its invasion of the Gaza Strip on Oct. 27 in retaliation to the Hamas terrorist attacks in Israel on Oct. 7.
More than 20,000 people have died in Gaza since Israel began its counterattack, according to the Hamas-controlled Gaza Health Ministry, and troops have destroyed roughly half of the buildings in the area.
But Sarraj faulted the Israeli military for something greater : the loss of the Gazan culture."The Israelis have also pulverized something else : Gaza City ' s cultural riches and municipal institutions,"Sarraj said."The unrelenting destruction of Gaza -- its @ @ @ @ @ @ @ @ @ @ archives and whatever economic prosperity it had -- has broken my heart."The op-ed criticized Israel for the destruction of many of the city ' s features, including its zoo, the main public library, the Children ' s Happiness Center, and more.
Sarraj accused the Israeli military of"destroying life"in Gaza, lamenting the country ' s"blockade of Gaza"that he says unfairly affects the Palestinian population."Why did the Israeli tanks destroy so many trees, electricity poles, cars and water mains?
Why would Israel hit a U.N. school?
The obliteration of our way of life in Gaza is indescribable.
I still feel I am in a nightmare because I can't imagine how any sane person could engage in such a horrific campaign of destruction and death,"he wrote.
The publication of the op-ed has sparked backlash from conservatives online who criticized the New York Times for amplifying the voice of a highly ranked Hamas member."Remember when the journalists lost their @ @ @ @ @ @ @ @ @ @ ran an op-ed by a sitting US Senator?
Well now the New York Times is running opinion pieces written by Hamas.
Are their journalists outraged?
NYT showing their true colors!"Chaya Raichik, creator of the popular Libs of TikTok page, wrote on X.
Raichik was referring to an instance in which a New York Times opinion editor resigned in summer 2020 after receiving backlash for publishing an op-ed from Sen.
Tom Cotton ( R-AR ) that called for using military force against violent participants in Black Lives Matter protests.
The editor, James Bennet, stepped down shortly after due to pushback not only from the public but also from employees inside the news organization.
Bennet later criticized the New York Times in a 17,000-word feature piece for the Economist, where he now works as a columnist and senior editor.
The former opinion editor accused the news organization of pushing a liberal bias, even going so far as to say the outlet censors differing opinions."I think many Times staff have little idea @ @ @ @ @ @ @ @ @ @ they are from fulfilling their compact with \textbf{readers} to show the world ' without fear or favor, '"he wrote."The Times ' s problem has metastasized from liberal bias to illiberal bias, from an inclination to favor one side of the national debate to an impulse to shut debate down altogether."

% matched lemmas: newsletter, reader
\end{textsample}
